Table~\ref{tab:physics} compares the cores
under the same SGT command sequence. Three observations stand out.
\emph{(i) Rigidity.} The DAE core keeps the bar lengths at
$\pDaeBar$\,m; the penalty+SHAKE core reaches $\pLegacyBar$\,m and, with
the interactive aids, $\nLegBar$\,m.
\emph{(ii) Hidden energy exchange.} To keep the bars rigid, the
penalty+SHAKE core removes $\pLegacyNum$\,J by position projection
while its cables do a comparable amount of work, i.e.\ the cables
mostly pump energy into the penalty bars that the projection then
discards. The DAE core changes only $\pDaeNum$\,J numerically. With the
original locomotion aids, \pLegacyAidShare\,\% of the energy input comes
from non-physical forces, which explains its high speed
(\nLegSpeed\,m/s) in Table~\ref{tab:nominal}.
\emph{(iii) Drift.} Before correction, both cores show a per-step energy
mismatch that is first-order discretization error: for seed~1, halving
$\Delta t$ from 2 to 1\,ms reduces the accumulated $\sum_k|\delta_k|$
from \dcDaeTwoAbs\ to \dcDaeOneAbs\,J for the DAE core and from
\dcLegTwoAbs\ to \dcLegOneAbs\,J for penalty+SHAKE, and the largest
single-step mismatch of the DAE core (at the initial landing) from
$\dcDaeTwoStep$ to $\dcDaeOneStep$\,J. The energy row removes this
drift every step, so $\varepsilon$ stays at the solver tolerance
($\pDaeResidual$\,J). Without the energy row the residual accumulates
to $\pNoERowResidual$\,J. We stress that the energy row corrects an
integration error of known order; it does not by itself prove that the
model is complete---that is the role of $\delta_k$, which would expose
an unbooked force as a non-convergent drift. On rugged Level~7 the
corrector occasionally reaches its 10-iteration cap during contact
transitions, leaving a transient residual of up to
$4\times10^{-2}$\,J (0.2\,\% of $E$) that is removed in the following
steps (Table~\ref{tab:terrain}).
